% =============================================================================
%  Part IV -- Units, parameters and the SI bridge (catalogue of symbols)
% =============================================================================

Causal Physics is formulated at the level of causal structure, not of spacetime geometry. Its
fundamental quantities are therefore dimensionless and relational: counts of events, relations
between them and ratios of such counts. Comparison with observation nevertheless requires
quantities reported in SI units. This part states how the two are kept apart, lists the symbols of
the theory, and states which numbers are parameters and which are bridges.

% =============================================================================
\section{Three descriptive layers}\label{sec:layers}
% =============================================================================

The theory separates three layers of description. The separation is a structural requirement, not
a convenience.

\begin{center}
\begin{tabularx}{\linewidth}{lXX}
\toprule
layer & role & permitted content\\
\midrule
causal layer & physical ontology and dynamics & events, the relation $\prec$, counts of events,
  flows, response variables, the action of the order\\
bridge layer & metrology only & the discreteness scale, unit conversion, projection of counts onto
  lengths, times and accelerations\\
observational layer & empirical comparison & SI quantities, catalogues, measured accelerations,
  delays, deflections\\
\bottomrule
\end{tabularx}
\end{center}

Three rules follow. Causal variables never carry physical dimensions. Observational quantities never
take part in causal evolution. Bridge parameters never change a rule, a threshold or a weight of the
causal layer. A construction that violates any of the three reintroduces background structure and
is excluded.

Coordinates have a special status. They are not part of the causal layer. An \emph{instrument} may
use them as a declared readout, for example to place the events of a sprinkling or to report a
profile as a function of radius; a \emph{rule} of the dynamics may not (Sec.~\ref{sec:method-ax}).

% =============================================================================
\section{Catalogue of symbols}\label{sec:symbols}
% =============================================================================

For Greek letters and special symbols the middle column gives the name and the pronunciation
(International Phonetic Alphabet, English).

\newcommand{\rd}[2]{#1~\textipa{/#2/}}
\newenvironment{symtab}{\begin{center}\small\begin{tabular}{l>{\raggedright\arraybackslash}p{4.0cm}>{\raggedright\arraybackslash}p{7.4cm}}
\toprule symbol & read as & meaning\\\midrule}{\bottomrule\end{tabular}\end{center}}

\subsection{Structure of the order}

\begin{symtab}
$v,\ u,\ x,\ y$ & & events (elements of the causal set)\\
$u\prec v$ & \rd{$u$ precedes $v$}{prI"si:dz} & causal relation: $u$ can influence $v$\\
$u\prec\!\!\cdot\, v$ & \rd{link}{lINk} & link: $u\prec v$ with nothing in between\\
$I(u,v)$ & & causal interval $\{w:u\prec w\prec v\}$\\
$n(u,v)$ & & number of elements strictly inside $I(u,v)$\\
$J^-(x)$ & \rd{J minus}{dZeI "maIn@s} & causal past of $x$\\
$\rho$ & \rd{rho}{roU} & \textbf{fundamental density of the substrate}: the mean number of events per unit
  proper four-volume. In a sprinkling the number of events in a region of volume $V$ is Poisson
  distributed with mean $\rho V$. It sets the discreteness length $\ell=\rho^{-1/D}$ and, with
  $\epsilon$, the non-locality scale $\ell_k$. It is the same everywhere: near matter the
  \emph{coordinate} density changes ($\rho_c$), not $\rho$\\
$D=d_s+1$ & & spacetime dimension; $d_s$ spatial dimensions\\
$N$ & & number of events in a region\\
\end{symtab}

\subsection{Transport, matter and action}

\begin{symtab}
$F_v$ & & flow carried by event $v$ (non-negative, conserved)\\
$\Delta Q_v$ & \rd{delta}{"dElt@} $Q$ & flow stored at $v$ (the residual of the balance at $v$)\\
$s_v$ & & source term at $v$\\
$u^a$ & & four-velocity carried by a flow as its state\\
$P^a=F\,u^a$ & & four-momentum of a flow; merging rule $F=|P|$\\
$T^{ab}$ & & energy--momentum tensor of flows, Eq.~\eqref{eq:Tab}\\
$e_0$ & & rest-frame energy density of a flow ($T_{ab}u^au^b$)\\
$C$ & & capacity: the largest flow an event can pass on\\
$\sigma_v=F_v/C$ & \rd{sigma}{"sIgm@} & load of event $v$\\
$\Delta\tau$ & \rd{delta tau}{"dElt@ taU} & proper time between successive events of a flow, counted by the order\\
$S_m=-\sum m\,\Delta\tau$ & $\Sigma$: \rd{sum}{s2m} & action of a flow (Axiom~\ref{ax:action})\\
$\mathcal R=\sum_v Z_v\,\Delta\tau$ & \rd{script R}{skrIpt A:r} & Causal Resistance Functional, the impedance-weighted action\\
\end{symtab}

\subsection{Response of the substrate}

\begin{symtab}
$\Pi_v$, $0<\Pi_v\le1$ & \rd{capital pi}{"k\ae pIt@l paI} & permeability, read as the local clock rate:
  $\dd\tau=\Pi\,\dd t$ along a static world line\\
$L=\ln\Pi$ & \rd{log pi}{lAg paI} & logarithmic clock potential; the stress coordinate of the substrate
  ($L<0$ near matter)\\
$I_v\ge1$ & & stiffness: resistance of the order to reconfiguration\\
$Z_v=I_v/\Pi_v\ge1$ & & impedance\\
$c_{\rm eff}$ & \rd{c effective}{si: I"fEktIv} & local speed of the causal cone; $c_{\rm eff}=c$ in the empty substrate\\
$\rho_c$ & \rd{rho c}{roU si:} & density of events per \emph{coordinate} volume, in units of $\rho$
  ($=\sqrt{-g}$ in a sprinkling); grows near matter\\
$\rho_s$ & \rd{rho s}{roU Es} & density of saturated events: the source of the static field in the rest
  frame of the matter ($T_{ab}u^au^b$ of the flows)\\
$q=\rho_c\,c_{\rm eff}$ & & causal throughput (conserved, Sec.~\ref{sec:response})\\
\end{symtab}

\subsection{Geometry read from the order}

\begin{symtab}
$\phi$ & \rd{phi}{faI} & clock potential, read from the number of elements of a small interval at rest\\
$\psi$ & \rd{psi}{saI} & spatial potential; $\phi+\psi$ is read from the boundary of the relation and
  $\phi-d_s\psi$ from the density\\
$\gamma=\psi/\phi$ & \rd{gamma}{"g\ae m@} & post-Newtonian parameter (light deflection, Shapiro delay)\\
$\beta$ & \rd{beta}{"beIt@} & post-Newtonian parameter of the non-linearity of the clock field\\
$U$ & & Newtonian potential ($U=-\phi$ in the weak field)\\
$\eta_{ab}$ & \rd{eta}{"eIt@} & Minkowski metric\\
$h_{ab},\ \bar h_{ab}$ & \rd{h bar}{eItS bA:r} & metric perturbation and its trace-reversed form\\
$\Box_g$ & \rd{box}{bAks}, d'Alembertian & wave operator of the metric $g$\\
$R$ & & scalar curvature\\
\end{symtab}

\subsection{Operators and the action of the order}

\begin{symtab}
$\epsilon$, $0<\epsilon\le1$ & \rd{epsilon}{"Eps@lAn} & smearing parameter of the Benincasa--Dowker operator\\
$\ell_k=(\epsilon\rho)^{-1/D}$ & \rd{ell k}{El keI} & non-locality scale\\
$C_i$ & & Benincasa--Dowker coefficients; in 3+1 $(-1,9,-16,8)$, fixed by $D$\\
$a_D$ & & normalisation; $a_4=4/\sqrt6$\\
$K_\epsilon(n)$ & \rd{K epsilon}{keI "Eps@lAn} & smeared kernel, Eq.~\eqref{eq:Bdiscrete}\\
$g_k(w)$ & & expected kernel, Eq.~\eqref{eq:gk}\\
$B_\epsilon$, $\Ex[B_\epsilon]$ & \rd{expectation of B}{Ek.spEk"teIS@n} & discrete BD operator and its expectation\\
$S_\BD$ & & BD action of the order, Eq.~\eqref{eq:SBD}\\
$V(x,y)$ & & proper volume of the interval between $y$ and $x$\\
$\hat K(k,T)$ & \rd{K hat}{keI h\ae t} & ratio of the second-order BD action to its Einstein--Hilbert value
  (Sec.~\ref{sec:bdaction}); $K_{\rm loc}$ its local part\\
$c_T$ & & coefficient of the age term of the BD action\\
$\cA,\ \cB,\ \cC$ & \rd{script A, B, C}{skrIpt eI bi: si:} & coefficients of the static quadratic form of the action of
  the order: conformal, mixed and optical sector (Einstein--Hilbert $6$, $8$, $2$), Sec.~\ref{sec:mixed}\\
$T_{\rm U}$ & & Unruh temperature (Sec.~\ref{sec:jacobson}); not the age $T$\\
$T$ & & age: how long a static field has existed in the past of an event\\
$\kappa$ & \rd{kappa}{"k\ae p@} & coupling of flows to the order\\
$\hbar$ & \rd{h-bar}{eItS bA:r} & reduced Planck constant\\
$\ell_P$ & \rd{ell P}{El pi:} & Planck length\\
\end{symtab}

% =============================================================================
\section{Parameters of the theory}\label{sec:parameters}
% =============================================================================

\paragraph{Fundamental.} The causal layer has two numbers that are not fixed by the order itself:
the density $\rho$, which sets the discreteness length $\ell=\rho^{-1/D}$, and the smearing
parameter $\epsilon$, which sets the non-locality scale $\ell_k=(\epsilon\rho)^{-1/D}\ge\ell$. All
results of Part~\ref{part:geometric} are stated in units of $\ell_k$; the expectation values of the
operator and of the action depend on $\epsilon$ and $\rho$ only through the product $\epsilon\rho$.
For $\epsilon=1$, $\ell_k=\ell$.

\paragraph{Fixed by the dimension.} The coefficients $C_i$ and the normalisation $a_D$ are not
parameters: they are fixed by requiring that the expected operator reproduce the d'Alembertian in
$D$ dimensions.

\paragraph{The coupling.} The strength with which a flow acts on the order, $\kappa$, becomes
Newton's constant in the bridge layer. It is not tuned in the causal layer: in the linear regime
it fixes only the normalisation of the field, never its shape or its tensor structure.

\paragraph{Working parameters of models.} Particular models of the response of the substrate (the
saturation level of a load-driven clock law, a materialisation threshold, a feedback exponent) carry
working parameters. They are properties of a model, not constants of nature, and no result of this
paper depends on their values.

\paragraph{Laboratory units.} The laboratory works in units in which $\rho=1$, or $\epsilon\rho=1$
for computations with the expected operator; lengths are then counted in $\ell$ or $\ell_k$.

% =============================================================================
\section{The SI bridge}\label{sec:bridge}
% =============================================================================

\paragraph{The speed of light.} In a Lorentz-invariant substrate the boundary of the relation has
the same slope in every frame. Reporting that slope in SI units defines $c$; $c$ is a conversion
between the counting of events along time and across space, not an additional constant of the
causal layer.

\paragraph{Newton's constant (conditional).} The bridge to $G$ goes through the entropy of horizons
(Sec.~\ref{sec:jacobson}). If the entropy of a horizon is the number of horizon molecules times an
entropy $s_0$ per molecule, the entropy density per unit area is $\eta=s_0\,a^{(4)}/\ell^2$ with
$a^{(4)}=\sqrt3/10$, and the thermodynamic derivation of Einstein's equations gives
$G=1/(4\hbar\eta)$~\cite{Jacobson1995}. Then
\begin{equation}\label{eq:lP}
  \ell^2=4s_0\,a^{(4)}\,\ell_P^2=\frac{2\sqrt3}{5}\,s_0\,\ell_P^2,\qquad
  \ell\simeq1.075\sqrt{s_0}\;\ell_P ,
\end{equation}
with $\ell_P$ the Planck length. The discreteness length is thereby tied to $G$ up to $s_0$. This is a
relation of consistency, not a prediction of $G$: it becomes predictive only when $\ell$ or $s_0$ is
measured independently.

\paragraph{The non-locality scale.} The smearing scale $\ell_k$ is not fixed by any bridge. It is a
free scale of the theory, to be bounded or measured by the predictions of Part~\ref{part:predictions}. For
fields propagated by the non-local operator existing data give $\ell_k\lesssim10^{-19}$\,m; whether this
applies to Causal Physics is discussed in Sec.~\ref{sec:window}.

\paragraph{Projection of counts.} Within the plateau, where geometry is a reliable compression of
the order, counts are projected onto SI quantities: a number of elements $N$ of a region onto a
four-volume $N/\rho$, the length of the longest chain onto a proper time, the clock potential $L$
onto the Newtonian potential $U/c^2$. Outside the plateau the projections persist formally but lose
physical meaning.

\paragraph{No acceleration scale.} The present formulation introduces no acceleration scale in any
layer (Sec.~\ref{sec:astro-cosmo}).
